\FloatBarrier
\noindent This study considers a bounded orbital-dynamics problem in near-polar low Earth orbit (LEO) intended to isolate the short-horizon interaction between solar radiation pressure (SRP), eclipse timing, atmospheric drag, and \(J_2\)-driven nodal precession. The emphasis is on whether a sail with inertially fixed orientation can exhibit short-horizon semi-major-axis growth while remaining close to Sun-synchronous orbital geometry.

\vspace{5pt}
\par\noindent A Sun-synchronous orbit (SSO) is desirable because it maintains a predictable Sun-relative geometry while allowing the orbital plane to precess with the apparent annual motion of the Sun \citep{vallado2007}. For near-circular LEO, this requires a secular right ascension of the ascending node (RAAN) precession rate of approximately
\begin{equation}
\frac{d\Omega}{dt} \approx 0.986^\circ/\text{day},
\label{eq:RAANrate}
\end{equation}
which is primarily supplied by Earth’s oblateness through the \(J_2\) perturbation. The corresponding nodal precession rate is \citep{vallado2007}
\begin{equation}
\dot{\Omega}
  = -\frac{3}{2}\,J_2\,
      \left(\frac{R_{E}}{a\bigl(1-e^{2}\bigr)}\right)^{\!2}
      \sqrt{\frac{\mu}{a^{3}}}\,
      \cos i.
\end{equation}
For circular orbits (\(e=0\)) in the altitude range considered here, the inclination required to satisfy Eq.~\eqref{eq:RAANrate} is close to \(98^\circ\), as summarised in Table~\ref{tab:sso_inclinations}.


\begin{table*}[htbp]
\centering
\caption{Required inclination to maintain Sun-synchronous orbits (SSO) for circular orbits (\(e = 0\)) at selected altitudes.}\label{tab:sso_inclinations}
\vspace{1em}
\begin{tabular}{c|c}
\textbf{Altitude (km)} & \textbf{Required Inclination (°)} \\
\hline
600 & 97.92 \\
700 & 98.04 \\
800 & 98.17 \\
900 & 98.31 \\
\end{tabular}

\end{table*}


\begin{table*}[htbp]
\centering
\caption{Baseline parameters used for the representative gain-case propagation and local mechanism analysis.}
\label{tab:baseline_params}
\vspace{0.5em}
\begin{tabular}{l c}
\hline
\textbf{Parameter} & \textbf{Value} \\
\hline
Initial altitude, \(h_0\) & \(700\,\mathrm{km}\) \\
Initial inclination, \(i_0\) & \(98.17^\circ\) \\
Initial eccentricity, \(e_0\) & \(1\times10^{-4}\) \\
Spacecraft mass, \(m\) & \(10\,\mathrm{kg}\) \\
Sail area, \(A\) & \(100\,\mathrm{m^2}\) \\
Area-to-mass ratio, \(A/m\) & \(10\,\mathrm{m^2\,kg^{-1}}\) \\
Reflectivity coefficient, \(C_r\) & \(1.7\) \\
Drag coefficient, \(C_D\) & \(2.2\) \\
\hline
\end{tabular}
\end{table*}

\vspace{5pt}
\par\noindent In the present work, the sail orientation is prescribed as inertially fixed over the time horizons of interest, and no active steering law is applied during propagation. This contrasts with actively steered sails \citep{dachwaldOptimalSolarSail2005,quartaOptimalInterplanetaryTrajectories2023} and with active momentum management approaches such as those developed for Solar Cruiser \citep{cannellaDesignOverviewSolar2021,shenSolarSailMomentum2026}. The fixed-attitude assumption is deliberate: it allows the orbital consequences of a fixed force geometry to be isolated directly, without introducing control-law dependence into the interpretation of the results. At the same time, it imposes an important limitation. As the orbital plane precesses and the Sun-relative geometry evolves, the orientation of the sail with respect to both the Sun and the local orbital frame cannot remain optimal indefinitely. Any favourable energy-transfer regime must therefore be understood as geometry-bounded and transient rather than as a permanently sustained operating mode.



\vspace{5pt}
\par\noindent The evolution of the orbit under weak non-Keplerian perturbations is interpreted using the Gauss Variational Equations (GVEs) \citep{vallado2007,battin1999}. Writing the perturbing acceleration vector \( \vec{F} \) in the local orbital frame as radial (\(F_R\)), transverse (\(F_T\)), and normal (\(F_N\)) components, the orbital-element evolution is given by
\begin{equation}
\frac{da}{dt} = \frac{2a^2}{h} \left[ e \sin{\theta} \, F_R + (1 + e \cos{\theta}) \, F_T \right],
\label{eq:da}
\end{equation}
\begin{equation}
\frac{de}{dt} = \frac{\mu}{h} \left[ \sin{\theta} \, F_R + \left( \cos{\theta} + \frac{e + \cos{\theta}}{1 + e \cos{\theta}} \right) F_T \right]
\end{equation}
\begin{equation}
\frac{di}{dt} = \frac{r \cos{u}}{h} \, F_N,
\label{eq:di}
\end{equation}
\begin{equation}
\frac{d\Omega}{dt} = \frac{r \sin{u}}{h \sin{i}} \, F_N,
\label{eq:domega}
\end{equation}
where \( h = \sqrt{\mu a (1 - e^2)} \) is the specific angular momentum, \( r = a(1 - e^2)/(1 + e \cos{\theta}) \) is the instantaneous orbital radius, \( \theta \) is the true anomaly, and \( u = \omega + \theta \) is the argument of latitude. A detailed derivation of the GVEs is given by \citet{vallado2007}.

\vspace{5pt}
\par\noindent These expressions clarify the role of each perturbing component. In near-circular LEO, semi-major-axis growth is driven primarily by the in-plane force, especially the transverse component \(F_T\), whereas inclination \(i\) and RAAN \(\Omega\) evolve through the normal component \(F_N\). Since Sun-synchronicity relies on the secular \(J_2\)-driven nodal precession described above, additional out-of-plane forcing should remain small if the orbit is to stay close to the desired SSO geometry.

\vspace{5pt}
\par\noindent This creates the central geometrical tension of the present study. To obtain semi-major-axis growth, the sail must produce a favourable in-plane force with a strong transverse projection, while remaining close to SSO requires the normal component \(F_N\) to stay limited. Because the sail orientation is fixed in inertial space, the Sun-relative cone angle and the orbital-frame force decomposition drift as the orbit evolves, so any favourable geometry is expected to persist only over a bounded interval.

\vspace{5pt}
\par\noindent % [PLACEHOLDER: Figure showing the time evolution of the SRP force decomposition in the orbital frame for a representative gain case. This should show \(F_T\) and \(F_N\) over time, or equivalently the Sun-relative geometric drift, to support the claim that favourable in-plane forcing is temporary.]
The force model retained for this study is intentionally simple and is chosen to isolate the mechanism of interest rather than to maximise environmental fidelity. The translational dynamics include central gravity, Earth’s oblateness through the \(J_2\) term, solar radiation pressure, eclipse shadowing, and atmospheric drag. No Earth albedo model is included.

\vspace{7pt}
\noindent\textbf{Solar radiation pressure} \vspace{2pt}

\par\noindent The SRP model is based on photon momentum transfer to a reflective flat plate \citep{mcinnes1999}. For an idealised square sail with cone angle \( \alpha \) relative to the Sun line, the force magnitude follows the familiar \( \cos^2\alpha \) dependence for a flat-plate sail (see \citealp{mcinnes1999}, Ch.~2). In the implementation used here, the SRP acceleration is written as
\begin{equation}
    \mathbf{a}_{\text{SRP}} = \frac{P_{1\ \text{AU}} \, A \, C_r}{m c} \cos^2{\alpha} \, \hat{\mathbf{r}},
    \label{eq:srp_force}
\end{equation}
where \(P_{1\ \text{AU}}\) denotes the solar radiation flux at one astronomical unit \citep{kopp2011solar}, \(A\) is the sail area, \(C_r\) is the reflectivity coefficient, \(m\) is the spacecraft mass, and \(\hat{\mathbf{r}}\) is the unit vector from the spacecraft to the Sun. The SRP acceleration is evaluated in the inertial frame and then resolved into the orbital frame for interpretation through the GVEs.

\vspace{5pt}
\par\noindent Eclipses are modelled using a binary conical shadow geometry. The spacecraft is treated as illuminated when it lies outside the Earth’s umbral cone and fully shadowed when it lies within it, so that SRP is either applied in full or suppressed entirely. This treatment is sufficient for the present purpose because the paper is concerned primarily with the orbital-phase selectivity of SRP application rather than with resolving penumbral transition details. The shadowed SRP acceleration is therefore written as
\begin{equation}
\mathbf{a}_{\text{SRP}}^{\text{shadowed}} = s \cdot \mathbf{a}_{\text{SRP}},
\end{equation}
where
\[
s =
\begin{cases}
1, & \text{illuminated},\\
0, & \text{eclipsed}.
\end{cases}
\]
This eclipse treatment is central to the present paper because the mechanism of interest depends not only on the magnitude of SRP, but on which orbital arcs remain illuminated and which are removed by shadow.

\vspace{5pt}
\par\noindent % [PLACEHOLDER: Geometry schematic showing the Sun direction, Earth shadow cone, orbit, and inertially fixed sail normal. This figure should visually identify illuminated and eclipsed arcs and should support the later discussion of selective removal of unfavourable SRP work arcs.]
\vspace{7pt}
\noindent\textbf{Atmospheric drag} \vspace{2pt}

\par\noindent Atmospheric drag is included because it remains dynamically relevant throughout the altitude range considered, particularly once eccentricity growth lowers perigee and increases local density exposure. The atmospheric density is represented using a piecewise-exponential fit to U.S.\ Standard Atmosphere (1976) data \citep{USSA1976} based on Braeunig’s tabulation \citep{braeunig_atmos}. Within each altitude interval \(i\), the density is evaluated as
\begin{equation}
\rho(h) = \rho_0^{(i)} \exp\left(-\frac{h - h_0^{(i)}}{H^{(i)}}\right),
\label{eq:atmospheric_density}
\end{equation}
where \(\rho_0^{(i)}\), \(h_0^{(i)}\), and \(H^{(i)}\) are the reference density, reference altitude, and scale height for that layer.

\vspace{5pt}
\par\noindent The aerodynamic model used here is a simple projected-area drag law. Atmospheric rotational velocity is not modelled; the relative flow is taken directly from the spacecraft inertial velocity. The drag acceleration is written as
\begin{equation}
\mathbf{a}_{\mathrm{drag}} =
-\frac{1}{2}\,\rho(h)\,v^{2}\,C_D\,\frac{A_{\mathrm{frontal}}}{m}\,\hat{\mathbf{v}},
\label{eq:drag_force}
\end{equation}
where \(C_D=2.2\) is a standard value for satellites in free-molecular flow \citep{vallado2007}, \(m\) is the spacecraft mass, \(v = \|\mathbf{v}\|\), and \(\hat{\mathbf{v}}\) is the unit vector in the velocity direction. The frontal area \(A_{\mathrm{frontal}}\) is computed from the projection of the sail surface onto the plane normal to the instantaneous velocity vector, so that the effective drag area varies with the angle between the inertially fixed sail normal and the flow.



\vspace{5pt}
\par\noindent The drag model is intentionally simple. Its role in this paper is not to resolve high-fidelity free-molecular aerodynamics, but to capture the dominant loss mechanism opposing SRP-driven energy gain. This is especially important once the orbit departs from circularity, because drag then remains concentrated near perigee and can overtake the orbit-averaged SRP surplus even when early-time semi-major-axis growth is observed.

\begin{table*}[htbp]
\centering
\caption{Summary of the principal modelling assumptions used in the orbital propagations.}
\label{tab:model_assumptions}
\vspace{0.5em}
\begin{tabular}{p{4.2cm} p{11.5cm}}
\hline
\textbf{Assumption / model element} & \textbf{Implementation used in this study} \\
\hline
Attitude treatment & Sail orientation prescribed as inertially fixed over the horizon of interest; no active steering law applied during propagation. \\

Central gravity & Included as the baseline point-mass Earth gravitational field. \\

\(J_2\) perturbation & Included to capture Earth oblateness and the associated secular nodal precession required for near-Sun-synchronous behaviour. \\

Solar radiation pressure (SRP) & Included using a flat-plate sail model with one-sided illumination and force scaling with incidence angle through a cosine-based projection law. \\

Eclipse model & Included using a conical Earth shadow model that suppresses SRP during eclipse intervals. \\

Atmospheric drag & Included using a projected-area flat-plate drag law with \(C_D = 2.2\). \\

Density model & Static piecewise-exponential atmospheric density fit based on U.S. Standard Atmosphere (1976) reference data. \\

Atmospheric co-rotation & Neglected; relative flow taken directly from the spacecraft inertial velocity. \\

Earth albedo & Omitted. \\

Orbit regime & Near-circular, near-polar LEO focused on maintaining approximate Sun-synchronous geometry. \\

Study objective & Controlled mechanism-isolation study of transient gain behaviour, not a full mission-feasibility or flight-operations model. \\
\hline
\end{tabular}
\end{table*}

\vspace{5pt}
\par\noindent % [PLACEHOLDER: Table summarising the model assumptions used in the propagation. This should include: inertially fixed sail orientation, binary conical eclipse model, SRP force law, projected-area drag law with \(C_D=2.2\), density model, omission of atmospheric co-rotation, and inclusion of \(J_2\).]
The assumptions adopted above are sufficient for the present purpose of mechanism isolation and regime mapping. They do not constitute a full validation of passive sail performance over long durations, nor do they resolve all structural, control, or environmental uncertainties relevant to flight operations. Those questions are outside the scope of this section and are addressed only insofar as they affect interpretation of the short-horizon transient dynamics studied in this paper.
\FloatBarrier